\begin{helper}
Let
\[
U_I=\noderef{TwoBiteUnclosedProjectedPairs}(\omega,I),\qquad
T_I=\noderef{TwoBiteTouchedProjectedPairs}(\omega,\varepsilon,I),
\]
and let
\[
O_I=\noderef{TwoBiteStagedOpenPairs}(\omega,\varepsilon,I).
\]
For any real number \(f\), if
\[
f-9\varepsilon_1|I|^2\le |U_I|
\qquad\text{and}\qquad
|T_I|\le \varepsilon_1|I|^2,
\]
then
\[
f-10\varepsilon_1|I|^2\le |O_I|.
\]
\end{helper}
\begin{proof}
All three sets are finite subsets of the same tagged projected-pair ambient
set
\[
P_I=\noderef{TwoBiteProjectionPairSet}(\omega,I).
\]
Write
\[
C(e)\Longleftrightarrow
\noderef{TwoBiteProjectionPairClosed}(\omega,I,e),
\]
\[
C_{\mathrm{sc}}(e)\Longleftrightarrow
\noderef{TwoBiteProjectionPairSameColorClosed}(\omega,I,e),
\]
\[
R(e)\Longleftrightarrow
\noderef{TwoBiteProjectionPairTouched}(\omega,\varepsilon,I,e),
\]
and
\[
Q(e)\Longleftrightarrow
e\in\noderef{TwoBitePreTerminalRecordedPairs}(\omega,\varepsilon,I).
\]
By the definitions,
\[
U_I=\{e\in P_I:\neg C(e)\},
\qquad
T_I=\{e\in P_I:R(e)\},
\]
and
\[
O_I=\{e\in P_I:\neg C_{\mathrm{sc}}(e)\wedge\neg R(e)\wedge\neg Q(e)\}.
\]

We first prove the inclusion
\[
U_I\setminus T_I\subseteq O_I. \tag{1}
\]
Fix \(e\in U_I\setminus T_I\).  Then \(e\in U_I\), so \(e\in P_I\) and
\(\neg C(e)\).  By
\(\noderef{TwoBiteSameColorClosedImpliesClosed}\), same-color closedness
implies closedness:
\[
C_{\mathrm{sc}}(e)\Longrightarrow C(e).
\]
Taking the contrapositive gives \(\neg C_{\mathrm{sc}}(e)\).  Since
\(e\notin T_I\) and \(e\in P_I\), the identity
\(T_I=\{e\in P_I:R(e)\}\) gives \(\neg R(e)\).

It remains to show \(\neg Q(e)\).  We prove the pointwise implication
\[
Q(e)\Longrightarrow C(e)\vee R(e). \tag{2}
\]
Unfold \(\noderef{TwoBitePreTerminalRecordedPairs}\).  If \(e\) is a red
tag, there are two possible branches.  In the incident-coordinate branch,
one red endpoint is stage-revealed and the other is projection-contained;
unfolding \(\noderef{TwoBiteProjectionPairTouched}\) shows that the red
tagged pair is touched, so \(R(e)\) holds.  In the \(X_x(I)\)-pair branch,
there is a stage-revealed base vertex \(x\) and the red coordinate lies in
\[
\noderef{TwoBiteBasePairs}
\bigl((\noderef{TwoBiteX}(\omega,I,x)).\mathrm{image}(\pi_R)\bigr).
\]
Unfolding \(\noderef{TwoBiteProjectionPairClosed}\) gives \(C(e)\), with
the same witness \(x\).  The blue-tagged case is identical, using
\(\pi_B\) in the \(X_x(I)\)-pair branch and the blue endpoints in the
incident-coordinate branch.  Hence (2) holds in every case.

Now if \(Q(e)\) held, (2) would give \(C(e)\vee R(e)\), contradicting the
already obtained \(\neg C(e)\) and \(\neg R(e)\).  Therefore \(\neg Q(e)\).
Together with \(e\in P_I\), \(\neg C_{\mathrm{sc}}(e)\), and \(\neg R(e)\),
this is exactly the defining membership condition for \(e\in O_I\).  This
proves (1).

Next compare cardinalities.  The set \(U_I\) is the disjoint union of its
part outside \(T_I\) and its part inside \(T_I\):
\[
U_I=(U_I\setminus T_I)\,\dot\cup\,(U_I\cap T_I).
\]
Therefore
\[
|U_I|=|U_I\setminus T_I|+|U_I\cap T_I|.
\]
Since \(U_I\cap T_I\subseteq T_I\),
\[
|U_I\cap T_I|\le |T_I|.
\]
It follows that
\[
|U_I|\le |U_I\setminus T_I|+|T_I|.
\]
After casting these natural cardinalities to real numbers and subtracting
\(|T_I|\), we get
\[
(|U_I|:\mathbb R)-(|T_I|:\mathbb R)
\le (|U_I\setminus T_I|:\mathbb R). \tag{3}
\]
By the inclusion (1), monotonicity of finite-set cardinality gives
\[
(|U_I\setminus T_I|:\mathbb R)\le (|O_I|:\mathbb R). \tag{4}
\]
Combining (3) and (4),
\[
(|U_I|:\mathbb R)-(|T_I|:\mathbb R)\le (|O_I|:\mathbb R). \tag{5}
\]

Now use the two numerical hypotheses.  The first says
\[
f-9\varepsilon_1|I|^2\le (|U_I|:\mathbb R),
\]
and the second says
\[
(|T_I|:\mathbb R)\le \varepsilon_1|I|^2.
\]
Subtracting the second inequality from the first yields
\[
f-9\varepsilon_1|I|^2-\varepsilon_1|I|^2
\le (|U_I|:\mathbb R)-(|T_I|:\mathbb R).
\]
Together with (5), this gives
\[
f-9\varepsilon_1|I|^2-\varepsilon_1|I|^2\le (|O_I|:\mathbb R).
\]
Finally,
\[
9\varepsilon_1|I|^2+\varepsilon_1|I|^2
=10\varepsilon_1|I|^2,
\]
so the last inequality is exactly
\[
f-10\varepsilon_1|I|^2\le |O_I|.
\]
\end{proof}
